\documentclass[10pt,a4paper]{amsart}
\usepackage[margin=19mm]{geometry}
\usepackage{amsmath,amssymb,mathtools}
\usepackage[scr=rsfs,cal=euler]{mathalfa}
\usepackage{xfrac}
\usepackage[unicode,hidelinks,bookmarks=false]{hyperref}
\newcommand{\ie}{i.e.\ }
\newcommand{\cf}{cf.\ }
\newcommand{\resp}{respectively\ }
\newcommand{\Subsection}{\subsection}
\providecommand{\nobreakdash}{\nobreak\hbox{-}}
\setlength{\emergencystretch}{2em}
\multlinegap0pt
\usepackage{bm}
\usepackage{bbm}
\usepackage{dsfont}
\usepackage{xspace}
\usepackage{cancel}
\usepackage{mathtools}
\usepackage{amsthm}
\makeatletter
\@ifundefined{Thm}{\newtheorem{Thm}{Theorem}}{}
\@ifundefined{Lem}{\newtheorem{Lem}[Thm]{Lemma}}{}
\makeatother
\newcommand{\Coker}{\operatorname{Cok}}
\newcommand{\Hom}{\operatorname{Hom}}
\newcommand{\Spec}{\operatorname{Spec}}
\newcommand{\bd}{\mathbf{d}}
\newcommand{\bil}[1]{\langle #1\rangle}
\newcommand{\cC}{{\mathcal C}}
\newcommand{\cO}{{\mathcal O}}
\newcommand{\calU}{\mathcal{U}}
\newcommand{\calZ}{\mathcal{Z}}
\newcommand{\dimv}{\underline{\dim}}
\newcommand{\irr}{\operatorname{Irr}}
\newcommand{\md}{\operatorname{mod}}
\newcommand{\rk}{\operatorname{rk}}
\title[Generically $\tau$-regular irreducible components of module ]{Generically $\tau$-regular irreducible components of module varieties}
\author{Grzegorz Bobi\'nski, Jan Schr\"oer}
\date{Source: arXiv:2502.13709v2 (CC BY 4.0)}
\begin{document}
\maketitle

\noindent\emph{Source and licence.} Generically $\tau$-regular irreducible components of module varieties. Version arXiv:2502.13709v2 (CC BY 4.0), distributed under the Creative Commons Attribution 4.0 International licence, as recorded in the arXiv licence metadata for this exact version and shown on the abstract page https://arxiv.org/abs/2502.13709v2. Full rights notice: licenses/arxiv-2502.13709-CC-BY-4.0.txt.\par\smallskip

\noindent\textit{Editorial scope.} Counted unit: the Thm whose source label is \texttt{thm:main0b} in the article (source0.tex lines 2712--2730, source label \texttt{thm:main0b}), delivered together with the article's own complete original proof of it (source0.tex lines 2732--2835, one \texttt{proof} environment of 104 lines). No mathematics is written, repaired or re-derived: statement and proof are the article's own text line for line. Required context retained uncounted: lem:gvector3 (lines 960--973); lem:presentationsemicont (lines 1274--1287); lem:maxrankfM (lines 1430--1447); lem:regular (lines 1654--1670); thm:bundle1 (lines 2318--2348); lem:bundle6 (lines 2560--2565). Every definition, display and label of those lines is retained unchanged. Omitted from this package and not redistributed: 1--959, 974--1273, 1288--1429, 1448--1653, 1671--2317, 2349--2559, 2566--2711, 2731--2731, 2836--4705. Nothing inside a retained excerpt is omitted or altered: every retained body is delivered exactly as the article's lines are written, so no omission receipt of any kind accompanies this package. Citations: the retained excerpts carry no citation command at all, so no bibliography entry of the article is transcribed and no reference is redistributed. Reference adaptation: none was needed and none was made; every reference inside the delivered unit resolves inside it, so no unresolved reference or citation is produced. Printed slips kept literal: no printed word, symbol, hypothesis or number of the counted statement or of its proof was altered. Layout and class: the article's own class is \texttt{amsart} with packages including amssymb, latexsym, graphicx, amscd, upgreek, lscape, xy, color, epic, eepic, xspace, mathrsfs, setspace, cases; the wrapper is the lane's pinned \texttt{amsart} editorial wrapper at 10pt on A4, which restates the theorem-like environments the retained text needs and adds the article's own macro definitions that the retained text needs (\texttt{\textbackslash Coker}, \texttt{\textbackslash Hom}, \texttt{\textbackslash Spec}, \texttt{\textbackslash bd}, \texttt{\textbackslash bil}, \texttt{\textbackslash cC}, \texttt{\textbackslash cO}, \texttt{\textbackslash calU}, \texttt{\textbackslash calZ}, \texttt{\textbackslash dimv}, \texttt{\textbackslash irr}, \texttt{\textbackslash md}, \texttt{\textbackslash rk}), with extra wrapper packages (bm, bbm, dsfont, xspace, cancel, mathtools). The delivered numbering is the unit's own. Assets: no retained span contains a figure environment, a graphics-inclusion command, a PDF-page inclusion, a table or a TikZ picture; no image file is copied into this package, the package declares no asset dependency and no image file is redistributed with it. Licence: version arXiv:2502.13709v2 of this article is distributed under the Creative Commons Attribution 4.0 International licence, as recorded in the arXiv licence metadata for this exact version and shown on the abstract page https://arxiv.org/abs/2502.13709v2. A full rights notice is delivered at licenses/arxiv-2502.13709-CC-BY-4.0.txt. Characters: every retained excerpt is pure ASCII, so the delivered engine, XeLaTeX with its default Latin Modern text font, reports no missing character.
\begin{Lem}
\label{lem:gvector3}
Let $M,N \in \md(A)$, and let
$$
P_1 \xrightarrow{f} P_0 \to M \to 0
$$
be a projective presentation.
Applying $\Hom_A(-,N)$ yields an exact sequence
\begin{multline*}
0 \to \Hom_A(M,N) \to \Hom_A(P_0,N) \to \Hom_A(P_1,N)
\\
\to \Hom_A(N,\tau_A(M)) \oplus \Hom_A(P_f,M) \to 0.
\end{multline*}
\end{Lem}

\begin{Lem}[{Fei}]\label{lem:presentationsemicont}
Let $\calZ$ be a closed irreducible subset of $\md(A,\bd)$.
Let $M \in \calZ$, and let
$$
P_1 \to P_0 \to M \to 0
$$ be
a projective presentation of $M$.
Then there is a dense open subset $\calU$ of $\calZ$ such that
$M \in \calU$ and
each $M' \in \calU$ has a projective presentation of the form
$$
P_1 \to P_0 \to M' \to 0.
$$
\end{Lem}

\begin{Lem}\label{lem:maxrankfM}
For $M \in \md(A)$ let
$$
P_1 \xrightarrow{f} P_0 \to M \to 0
$$
be a projective presentation.
Then the following are equivalent:
\begin{itemize}\itemsep2mm

\item[(i)]
$\rk(f) = r(P_1,P_0)$;

\item[(ii)]
$\rk(f_M) = r(P_1^M,P_0^M)$ and
$\Hom_A(P_f,M) = 0$.

\end{itemize}
\end{Lem}

\begin{Lem}\label{lem:regular}
For
$M \in \md(A,\bd)$ the following are equivalent:
\begin{itemize}\itemsep2mm

\item[(i)]
$c(M) = e(M)$;

\item[(ii)]
$M$ is a regular point of the affine scheme
$\Spec(A,\bd)$.

\end{itemize}
In this case,
$M$ is contained in exactly one irreducible component
of $\md(A,\bd)$.
\end{Lem}

\begin{Thm}\label{thm:bundle1}
The following hold:
\begin{itemize}\itemsep2mm

\item[(i)]
Let $\cC$ be a closed irreducible stable subset of
$\Hom_A(P_1,P_0,\bd)$, and let $\calZ := \eta_\bd(\cC)$.
Then
$$
c(\calZ) = E(\calZ) + \hom_A(P_f,M) - {\rm codim}(\cC)
$$
where $f$ is generic in $\cC$ and $M$ is generic in $\calZ$, and
$$
{\rm codim}(\cC) := \hom_A(P_1,P_0) - \dim(\cC).
$$
Furthermore,
$\hom_A(P_f,M) - {\rm codim}(\cC) \le 0$.

\item[(ii)]
If $\cC := \Hom_A(P_1,P_0,\bd)$ is irreducible, then
$\calZ := \eta_\bd(\cC) \in \irr(A,\bd)$.

\item[(iii)]
The stratum
$\cC := \Hom_A(P_1,P_0)^\circ$ is irreducible and
$$
\calZ := \eta_\bd(\cC) \in \irr^\tau(A,\bd_{P_1,P_0}^\circ).
$$

\end{itemize}
\end{Thm}

\begin{Lem}\label{lem:bundle6}
Assume that $\cC := \Hom_A(P_1,P_0,\bd)$ is an irreducible stratum of
$\Hom_A(P_1,P_0)$, and let
$\calZ := \eta_\bd(\cC)$.
Then for each $f \in \cC$ we have $\Coker(f) \in \calZ^{\rm int}$.
\end{Lem}

\begin{Thm}\label{thm:main0b}
For $M \in \md(A)$ let
\begin{equation}\label{eq:presentation}
P_1 \xrightarrow{f} P_0 \to M \to 0
\end{equation}
be a projective presentation.
Then
the following hold:
\begin{itemize}\itemsep2mm

\item[(i)]
If $\rk(f) = r(P_1,P_0)$, then $M$ is $\tau$-regular;

\item[(ii)]
If $M$ is $\tau$-regular and if (\ref{eq:presentation}) is a minimal presentation,
then $\rk(f) = r(P_1,P_0)$.

\end{itemize}
\end{Thm}

\begin{proof}
Let $\bd := \dimv(M)$.

(i):
Let $\calZ := \calZ_{P_1,P_0}$.
We assume that $\rk(f) = r(P_1,P_0)$,
i.e.\ $\bd = \bd_{P_1,P_0}^\circ$.
Thus $\cC := \Hom_A(P_1,P_0,\bd) = \Hom_A(P_1,P_0)^\circ$ is irreducible,
$\calZ = \eta_\bd(\cC)$,
and Lemma~\ref{lem:bundle6} implies that
$M \in \calZ^{\rm int}$.
Since $\rk(f) = r(P_1,P_0)$, Lemma~\ref{lem:maxrankfM}
says that $\Hom_A(P_f,M) = 0$.
By Lemma~\ref{lem:gvector3}
there is an exact sequence
$$
0 \to \Hom_A(M,M) \to \Hom_A(P_0,M) \to \Hom_A(P_1,M)
\to \Hom_A(M,\tau_A(M)) \to 0.
$$
The alternating sum of the dimensions of these
four spaces is zero.
We have
$$
c(M) = \dim(\calZ) - \dim \cO_M = \dim(\calZ) - \dim G_\bd + \hom_A(M,M).
$$
Recall that $c(M) \le E(M) = \hom_A(M,\tau_A(M))$.

Now an easy calculation yields that
$c(M) = E(M)$ if and only if
$$
\dim(\calZ) - \dim G_\bd = -\hom_A(P_0,M) + \hom_A(P_1,M).
$$
Note that
these numbers only depend on $\dimv(M)$.

Since $\calZ$ is generically $\tau$-regular, there exists some
$N \in \calZ$ such that $c(N) = E(N)$ and $N$ has
a projective presentation
$$
P_1 \xrightarrow{g} P_0 \to N \to 0.
$$
(For the latter we used Fei's Lemma~\ref{lem:presentationsemicont}.)
Again by Lemma~\ref{lem:gvector3}
we get an exact sequence
\begin{multline*}
0 \to \Hom_A(N,N) \to \Hom_A(P_0,N) \to \Hom_A(P_1,N)
\\
\to \Hom_A(N,\tau_A(N)) \oplus \Hom_A(P_g,N) \to 0.
\end{multline*}
Of course, we have $\dimv(M) = \dimv(N)$ and therefore
$\rk(g) = r(P_1,P_0)$.
This implies $\Hom_A(P_g,N) = 0$, see again Lemma~\ref{lem:maxrankfM}.
Now $c(N) = E(N)$ together with the considerations above implies that
$c(M) = E(M)$.


(ii):
Assume that $M$ is $\tau$-regular, i.e.\ $c(M) = E(M)$
and that the presentation (\ref{eq:presentation}) is minimal.
By Lemma~\ref{lem:regular}
there is a unique $\calZ \in \irr(A,\bd)$ with $M \in \calZ$, and by
upper semicontinuity, $\calZ$ is generically $\tau$-regular.
For $N$ generic in $\calZ$ we have $c(N) = E(N)$ and there is a projective presentation
$$
P_1 \xrightarrow{g} P_0 \to N \to 0.
$$
(Here we used again Lemma~\ref{lem:presentationsemicont}.)
We have
$$
g(N) = [P_1] - [P_0] - [P_g].
$$
For all $L \in \calZ^{\rm int}$ we have
\begin{align*}
\bil{g(L),[L]} &= - \hom_A(L,L) + \hom_A(L,\tau_A(L))
\\
&= E(L) - c(L) + \dim(\calZ) - \dim G_\bd.
\end{align*}
Now $c(M) = E(M)$ and $c(N) = E(N)$ imply that
$$
\bil{g(M),[M]} = \dim(\calZ) - \dim G_\bd = \bil{g(N),[N]}.
$$
We get
$$
\bil{[P_1]-[P_0],[M]} = \bil{g(M),[M]} = \bil{g(N),[N]} =
\bil{[P_1]-[P_0],[N]} - \bil{[P_g],[N]}.
$$
Since $[M] = [N] = \bd$, this
implies $\hom_A(P_g,N) = \bil{[P_g],[N]} = 0$.

Let $\cC \in \irr(\Hom_A(P_1,P_0,\bd))$ such that
$\eta_\bd(\cC) = \calZ$.
We can assume that $g$ is generic in $\cC$ and $N$ is generic in
$\calZ$.
Now Theorem~\ref{thm:bundle1}(i) says that
$$
c(\calZ) = E(\calZ) + \hom_A(P_g,N) - {\rm codim}(\cC)
= E(\calZ) - {\rm codim}(\cC).
$$
Since $c(\calZ) = E(\calZ)$,
this implies that ${\rm codim}(\cC) = 0$, and therefore
$\bd = \bd_{P_1,P_0}^\circ$.
In other words, $\rk(g) = r(P_1,P_0)$.
Since $\dimv(M) = \dimv(N)$ we also get $\rk(f) = r(P_1,P_0)$.
\end{proof}

\end{document}
